% Living methods document for `demiurge`, per the refactor charter
% (HANDOFF3 Sec. 7): one section per additive channel, mirroring
% `src/demiurge/`'s own module layout, updated in the SAME COMMIT as any
% change to a citation-bearing functional form or Tier 2 default prior in
% the code -- not a separate write-up pass done later.
%
% Target: AAS Journals (ApJ), AASTeX v7.0.1 (current as of 2026-09,
% https://journals.aas.org/aastexguide/). Compile with a TeX distribution
% that has aastex7.cls available (current TeX Live/MiKTeX, or Overleaf's
% built-in AAS Journals template) -- the class file is not vendored in this
% repo.
%
% Deliberately mostly empty right now: no demiurge channel modules have
% been built yet (see src/demiurge/parameters/registry.py's own module
% docstring -- same "scaffolding now, populate per real feature" discipline
% applies here). Add a \subsection per channel alongside that channel's
% real code, not ahead of it.

\documentclass{aastex7}

\shorttitle{demiurge}
\shortauthors{TBD}

\begin{document}

\title{demiurge: TBD Title}

% TODO: real author list/affiliations/emails once decided -- AASTeX v7
% requires \email for every author.
\author{TBD}
\affiliation{TBD}
\email{TBD}

\begin{abstract}
TBD -- written once the methods below are substantially populated, not
before.
\end{abstract}

% TODO: Unified Astronomy Thesaurus (UAT) concept keywords, once relevant
% (https://astrothesaurus.org/).
% \keywords{\uat{Concept}{ID}}

\section{Introduction}
TBD.

\section{Methods}
Each subsection below corresponds to one additive channel in
\texttt{src/demiurge/} and is added only once that channel's real code
exists, per this document's own header note.

\subsection{Galaxy Continuum} \label{sec:galaxy-continuum}
Mirrors \texttt{src/demiurge/galaxy\_continuum/}. Produces the pure stellar
continuum for the galaxy mock path -- the \texttt{continuum} bucket's main
content. Dust attenuation is explicitly excluded from this channel
(\texttt{dust1}=\texttt{dust2}=0 in every FSPS call); it is the
responsibility of a separate, not-yet-built channel, following the 3-bucket
decomposition's own scope-separation convention.

\subsubsection{Star Formation History} \label{sec:sfh}
The star-formation history SFR($t$) is generated non-parametrically,
following the Dense Basis method of \citet{IyerGawiser2017} and
\citet{Iyer2019}. Three stellar-mass-formation-time quantiles
($t_{25}$, $t_{50}$, $t_{75}$, i.e.\ the lookback times at which 25\%,
50\%, and 75\% of the galaxy's total stellar mass has formed) are drawn
from a Dirichlet prior over the four time-gaps they divide $[0, t_{\rm
obs}]$ into (symmetric, concentration $\alpha=2$; inspired by, but not
independently verified against, the $\alpha=1$ per-bin sSFR-fraction prior
of \citet{IglesiasNavarro2024}, who apply a closely related Dense-Basis-
based method to simulation-based inference of SFHs from DESI-like spectra).
A Gaussian Process with a Mat\'ern-3/2 kernel is then conditioned to pass
exactly through these constraints (plus the fixed endpoints $F(0)=0$,
$F(t_{\rm obs})=1$) in cumulative-stellar-mass-fraction-vs-time space,
following \citet{Iyer2019}'s Section 2.1 method; SFR($t$) is obtained by
numerically differentiating the reconstructed curve, clipped non-negative
and rescaled to conserve the drawn total stellar mass exactly.

Two implementation departures from \citet{Iyer2019}, adopted for simplicity
(\texttt{sfh.py}'s module docstring): the reported stochasticity is the
Gaussian Process posterior \emph{mean} through randomly-drawn quantile
constraints, not a draw from the full posterior; and the smoothness/
burstiness length-scale hyperparameter's default range is not independently
verified against \citeauthor{Iyer2019}'s own calibrated value.

\subsubsection{Metallicity History} \label{sec:metallicity}
Rather than an independent free function of time, the metallicity history
$Z(t)$ is tied directly to the star formation history via the classical
closed-box chemical evolution relation of \citet{SearleSargent1972}:
\begin{equation}
    Z(t) = y \, \ln\left(\frac{1}{\mu(t)}\right), \qquad
    \mu(t) = 1 - \epsilon \, F(t),
\end{equation}
where $y$ is an effective nucleosynthetic yield, $\epsilon$ is a star-
formation efficiency (the fraction of the initial gas reservoir ultimately
locked into stars by $t_{\rm obs}$), and $F(t)$ is the same cumulative
stellar-mass fraction from Section~\ref{sec:sfh}. $Z(t)$ is monotonically
non-decreasing by construction. This is a deliberate first-pass
simplification: a more complete inflow/outflow-aware, SFH-regulated
chemical evolution model exists \citep{YinShenHao2023} but requires
machinery (a Kennicutt--Schmidt law, mass-loading-factor-parameterized
outflows) well beyond this pass's scope.

\subsubsection{Initial Mass Function} \label{sec:imf}
The full theory of a time-varying, galaxy-wide stellar initial mass
function (the Integrated Galaxy-wide IMF, IGIMF: \citealt{KroupaWeidner2003};
modern formulation \citealt{Jerabkova2018}) derives the galaxy-wide IMF by
integrating a per-embedded-cluster stellar IMF against an embedded-cluster
mass function whose slope depends on the instantaneous star-formation rate,
with the maximum stellar mass per cluster set by a transcendental
mass-normalization condition. This is deliberately \emph{not} implemented
here, in favor of a much simpler treatment prioritizing a working,
observationally-testable channel over complete IMF physics.

What \emph{is} implemented is the metallicity-dependent shift of the
canonical \citet{Kroupa2001} IMF's low-mass (0.08--0.5~$M_\odot$) and
intermediate-mass (0.5--1~$M_\odot$) power-law slopes,
\begin{equation}
    \alpha_{1,2}(Z) = \alpha_{1,2}^{\rm canonical} + 0.5\,[M/H],
\end{equation}
with $[M/H] = \log_{10}(Z/Z_\odot)$ (using $Z_\odot = 0.0142$,
\citealt{Asplund2009}, matching the MIST isochrones used throughout),
$\alpha_1^{\rm canonical}=1.3$, $\alpha_2^{\rm canonical}=2.3$
\citep{Kroupa2001}. This relation was verified directly against the source
code of \texttt{galIMF} \citep{YanJerabkovaKroupa2017,
YanJerabkovaKroupaVazdekis2019}, the public reference implementation of the
IGIMF framework, rather than inferred from the papers' text alone. The
high-mass slope $\alpha_3$ is held fixed at its canonical value (2.3)
regardless of $Z$ or SFR in this implementation -- the SFR-driven,
cluster-mass-function-integration-dependent behavior IGIMF theory predicts
is not captured. A fuller physical treatment, deriving the IMF's
characteristic mass from gas thermodynamics (dust-gas coupling, metal-line
cooling, radiative feedback) rather than an empirical SFR/$Z$ scaling
\citep{ShardaKrumholz2022}, and potentially incorporating AGN radiative
feedback as an additional regulating mechanism (no literature precedent
identified for the latter -- would be a genuine extension of this project's
own, not a citation), is deferred future work.

\subsubsection{Spectral Synthesis}
Given SFR($t$), $Z(t)$, and (in the ``dynamic'' IMF mode) $\alpha_{1,2}(Z(t))$
from the sections above, the composite stellar spectrum is synthesized via
the Flexible Stellar Population Synthesis code \citep[FSPS;][]{Conroy2009,
Conroy2010}, using the MIST isochrones and the high-resolution C3K
spectral library. Two numerically distinct implementations are provided,
verified to agree (root-mean-square relative flux deviation
$\lesssim$ a few percent on held-out draws; see \texttt{demiurge}'s test
suite):

\begin{itemize}
    \item \textbf{fsps\_direct}: the star-formation history is split into
    several contiguous time segments, each assigned its own
    SFR-weighted mean metallicity (and, in dynamic mode, IMF slopes); FSPS
    is called once per segment via its tabulated-SFH mode, and the
    resulting spectra are summed (valid since stellar population spectra
    are linear in formed mass at fixed age, metallicity, and IMF). This is
    the slower (order 10--25~s per FSPS call) but always-available
    implementation, used as the validation reference for the alternative
    below.
    \item \textbf{pretabulated} (default): exploits the empirical finding
    that FSPS caches its full native age-grid response for a given
    (metallicity, IMF) setting on first evaluation, making subsequent
    age queries at that same setting orders of magnitude faster than a new
    setting. A modest, log-spaced grid of FSPS's full age response at
    fixed metallicity values (log-spaced: empirically substantially more
    accurate per grid point than linear spacing, since realized $Z(t)$
    histories spend most of their dynamic range at low $Z$) is
    precomputed once, offline, and shipped as static data; spectra are
    then reconstructed via bilinear interpolation (in log-age and log-$Z$)
    of that grid, implemented in both NumPy and PyTorch (auto-detecting
    CUDA), with no further calls to FSPS. This reduces per-object synthesis
    time from order 10~s to order 1--10~ms.
\end{itemize}

\section{Discussion}
TBD.

\bibliography{refs}
\bibliographystyle{aasjournal}

\end{document}
